\chapter{Real-time Reorienting During Episodic Retrieval Rescues Memories from Suboptimal Arousal-Attentional States}

Moment-to-moment fluctuations in preparatory attention and arousal explain, in part, why certain moments are more conducive to successful memory than others [1–14]. Evidence from neuroimaging, electrophysiology, and pupillometry indicates that such fluctuations reflect systematic shifts in the brain’s arousal and attentional states — henceforth, ‘A-A states’ — including states modulated by the locus coeruleus-norepinephrine (LC-NE) system [15–23]. Correlations between read-outs of A-A states, cognitive control, and memory motivate readiness-to-learn and readiness-to-remember frameworks [5,6,14,24], which build on a rich literature linking preparatory states to cognitive performance [25–39]. Notably, the apparent effects of A-A states at the time of retrieval highlight fundamental challenges for education, clinical populations with attention dysregulation, and everyday cognitive performance: well-encoded memories are not sufficient for performance if attention and arousal fluctuations in the moment just prior to an attempt to remember decrease the probability of successful memory retrieval.
